\chapter{Introducción}
\label{cap:introduccion}

\section{Contexto y motivación}

La detección de intrusiones en entornos corporativos ha evolucionado en las últimas dos
décadas desde el análisis de firmas sobre tráfico de red hasta plataformas que agregan
telemetría de múltiples fuentes: cortafuegos, sistemas de detección de red, agentes de
\textit{endpoint}, servicios en la nube y registros de aplicación. Esa evolución ha
resuelto un problema de \textit{visibilidad}, pero ha creado otro de \textit{escala}.

Un centro de operaciones de seguridad (SOC) de tamaño medio recibe diariamente un volumen
de alertas que excede con claridad su capacidad de análisis. La consecuencia directa es la
llamada \textbf{fatiga de alertas}: cuando la proporción de avisos irrelevantes es
suficientemente alta, el analista deja de tratarlos individualmente y adopta estrategias de
descarte que, inevitablemente, acaban filtrando también señales legítimas. El problema no
es únicamente de cantidad, sino de \textbf{ausencia de contexto}. Una alerta que informa de
que un proceso ha accedido a la memoria de otro no dice, por sí sola, si ese comportamiento
corresponde a una técnica adversaria documentada, si ha sido observada previamente en
campañas de un actor conocido, ni qué activos de la organización quedan comprometidos si la
sospecha se confirma.

A esta situación se añade una segunda dificultad, de naturaleza estructural. Los ataques
dirigidos no se manifiestan como eventos aislados sino como \textbf{secuencias de etapas}
encadenadas en el tiempo. Un adversario que busca credenciales privilegiadas ejecuta primero
un reconocimiento del entorno, obtiene después acceso a un equipo, carga las herramientas
necesarias, extrae la información y, con frecuencia, borra sus huellas. Cada una de esas
acciones, considerada por separado, puede resultar indistinguible de la actividad
administrativa legítima. Es la \textit{secuencia}, y no cada evento individual, la que
constituye la evidencia sólida del ataque. Los sistemas de detección que evalúan eventos de
forma independiente son, por construcción, incapaces de reconocer esa estructura.

El panorama de amenazas actual agrava ambos problemas. Los informes de referencia del sector
señalan un crecimiento sostenido de los ataques asistidos por inteligencia artificial y una
consolidación de MITRE ATT\&CK \cite{mitre_attack} como lenguaje común para describir el
comportamiento adversario \cite{enisa2025}. Existe por tanto una necesidad clara de
\textbf{pasar de la alerta aislada al incidente correlacionado, contextualizado y
priorizado}, que es precisamente el ámbito de la denominada \textit{conciencia
cibersituacional}.

\section{Planteamiento del problema}

El problema que aborda este trabajo puede formularse en tres preguntas encadenadas:

\begin{enumerate}
    \item ¿Cómo detectar automáticamente que un conjunto de eventos dispersos, procedentes
    de fuentes heterogéneas, constituye en realidad una \textbf{secuencia de ataque
    coherente} y no actividad legítima?
    \item Una vez detectada, ¿cómo determinar el \textbf{alcance real del riesgo} que esa
    secuencia supone para la organización, en función de los activos que puede llegar a
    comprometer?
    \item ¿Cómo presentar el resultado de forma que el analista pueda \textbf{confiar en la
    priorización y comprender su fundamento}, en lugar de recibir una puntuación opaca?
\end{enumerate}

Las tres preguntas comparten un requisito transversal: la respuesta debe producirse en
tiempo útil sobre un flujo continuo de datos, no como resultado de un análisis
\textit{a posteriori}.

\section{Objetivos}

\subsection{Objetivo general}

Diseñar, implementar y validar una plataforma de datos que detecte ataques multi-etapa en
tiempo real y priorice los incidentes resultantes de forma explicable, integrando
inteligencia de amenazas y razonamiento sobre rutas de ataque para dar soporte a la
conciencia cibersituacional.

\subsection{Objetivos específicos}

\begin{enumerate}
    \item \textbf{Ingerir y normalizar} telemetría heterogénea de red y de sistema operativo
    a un esquema canónico común, sobre una arquitectura de procesamiento en
    \textit{streaming}.

    \item \textbf{Detectar secuencias de ataque multi-etapa} mediante procesamiento de
    eventos complejos, alineando las detecciones con las tácticas y técnicas de MITRE
    ATT\&CK.

    \item \textbf{Construir y mantener un grafo de ataque} que integre activos, topología,
    vulnerabilidades e inteligencia de amenazas estructurada.

    \item \textbf{Aplicar razonamiento lógico} sobre dicho grafo para derivar las rutas de
    ataque factibles hacia los activos críticos de la organización.

    \item \textbf{Priorizar los incidentes de forma explicable}, combinando la confianza de
    la detección, la severidad de la etapa alcanzada, la criticidad del activo afectado y su
    proximidad en una ruta de ataque derivada.

    \item \textbf{Validar cuantitativamente} el sistema sobre conjuntos de datos etiquetados
    con MITRE ATT\&CK, y proporcionar un despliegue reproducible.
\end{enumerate}

\section{Metodología de trabajo}

El desarrollo sigue un enfoque \textbf{incremental por fases}, en el que cada fase produce
un sistema funcional y verificable antes de acometer la siguiente. Esta decisión responde a
dos consideraciones prácticas: por un lado, garantiza que en cualquier momento del proyecto
exista una versión demostrable; por otro, permite validar las hipótesis técnicas de mayor
riesgo de forma temprana, antes de invertir esfuerzo sobre una arquitectura que pudiera
resultar inviable.

Cada componente se valida sobre datos reales de ataques documentados desde el momento de su
construcción, en lugar de posponer la verificación a una fase final de pruebas. Este
principio ha demostrado su utilidad durante el desarrollo: varios de los problemas más
relevantes descritos en el Capítulo~\ref{cap:desarrollo} son errores silenciosos que no
producían fallo alguno y que solo se detectaron al inspeccionar la coherencia de los
resultados intermedios.

Todas las decisiones de arquitectura relevantes se han registrado de forma continua junto
con su justificación y las alternativas descartadas, práctica que constituye la base del
apartado de decisiones de diseño de este documento.

\section{Estructura del documento}

El resto de la memoria se organiza como sigue:

\begin{itemize}
    \item El \textbf{Capítulo~\ref{cap:estadoarte}} revisa el estado del arte: los modelos y
    estándares sobre los que se apoya el trabajo (ATT\&CK, OCSF, STIX, Sigma), las técnicas
    de detección de secuencias, los grafos de ataque y el razonamiento lógico aplicado a
    seguridad, y los trabajos relacionados, identificando el hueco que este trabajo pretende
    cubrir.

    \item El \textbf{Capítulo~\ref{cap:desarrollo}} describe el análisis de requisitos, la
    arquitectura propuesta, las decisiones de diseño adoptadas y la implementación de cada
    uno de los componentes del sistema.

    \item El \textbf{Capítulo~\ref{cap:resultados}} presenta el protocolo de evaluación, los
    experimentos realizados y el análisis de los resultados obtenidos frente a las líneas
    base establecidas.

    \item El \textbf{Capítulo~\ref{cap:conclusiones}} recoge las conclusiones del trabajo y
    las líneas de continuación identificadas.

    \item Los capítulos finales abordan los aspectos éticos, sociales y legales del trabajo,
    y su valoración económica.
\end{itemize}
